\section{Swim in Rising Water}
\label{sec:graph-swim-rising-water}

\problemstatement{Given an $n\times n$ grid where
$\textit{grid}[r][c]$ is the elevation of a cell, at time $t$ the water
level is $t$ and you may move (in four directions) into any adjacent
cell whose elevation is $\le t$. Starting at $(0,0)$, find the least
time to reach $(n-1,n-1)$.}

\begin{patternbox}
\emph{``Least time until a path exists,'' where time is a rising
threshold}, is a connectivity-over-time question, and
\textbf{Disjoint Set} answers it beautifully: activate cells in
increasing elevation order, union each with its already-active
neighbors, and the answer is the elevation at which the start and end
first become connected.
\end{patternbox}

\subsection*{Turning a threshold into unions over time}

At time $t$, exactly the cells with elevation $\le t$ are usable. As $t$
rises, cells switch on one by one in increasing-elevation order. When a
cell switches on, it may bridge previously separate open regions -- a
$\textit{union}$ with each adjacent already-active cell. The moment
$(0,0)$ and $(n-1,n-1)$ land in the same set, a fully-submerged-enough
path exists, and the elevation of the just-activated cell is the minimum
time needed.

\begin{ideabox}[Process Cells in Elevation Order]
Sorting cells by elevation and activating them in that order means the
current cell's elevation is the current water level $t$. Because we stop
the instant start and end connect, that elevation is exactly the least
$t$ for which a path exists -- equivalently, the minimum over all paths
of the maximum elevation on the path.
\end{ideabox}

\subsection*{Algorithm}

\begin{enumerate}
  \item Collect all cells as $(\textit{elevation}, r, c)$ and sort
        ascending.
  \item Create a Disjoint Set over the $n^2$ cells; keep an
        \textit{active} flag per cell.
  \item Activate cells in order. For each activated cell, unite it with
        every adjacent active cell. After each activation, if $(0,0)$
        and $(n-1,n-1)$ are connected, return the current cell's
        elevation.
\end{enumerate}

\subsection*{C++ implementation}

\begin{lstlisting}
#include <bits/stdc++.h>
using namespace std;

class DisjointSet {
    vector<int> parent, size;
public:
    DisjointSet(int n) : parent(n), size(n, 1) {
        for (int i = 0; i < n; i++) parent[i] = i;
    }
    int find(int x) {
        if (parent[x] == x) return x;
        return parent[x] = find(parent[x]);
    }
    void unite(int x, int y) {
        int rx = find(x), ry = find(y);
        if (rx == ry) return;
        if (size[rx] < size[ry]) swap(rx, ry);
        parent[ry] = rx; size[rx] += size[ry];
    }
    bool connected(int x, int y) { return find(x) == find(y); }
};

int swimInWater(vector<vector<int>>& grid) {
    int n = grid.size();
    vector<tuple<int,int,int>> cells;              // {elevation, r, c}
    for (int r = 0; r < n; r++)
        for (int c = 0; c < n; c++)
            cells.push_back({grid[r][c], r, c});
    sort(cells.begin(), cells.end());

    DisjointSet ds(n * n);
    vector<vector<bool>> active(n, vector<bool>(n, false));
    int dr[] = {-1,1,0,0}, dc[] = {0,0,-1,1};

    for (auto& [elev, r, c] : cells) {
        active[r][c] = true;
        for (int k = 0; k < 4; k++) {
            int nr = r + dr[k], nc = c + dc[k];
            if (nr < 0 || nr >= n || nc < 0 || nc >= n) continue;
            if (active[nr][nc]) ds.unite(r * n + c, nr * n + nc);
        }
        if (ds.connected(0, n * n - 1)) return elev;   // start & end joined
    }
    return -1;
}

int main() {
    vector<vector<int>> grid = {{0,2},{1,3}};
    cout << swimInWater(grid) << endl;             // 3
    return 0;
}
\end{lstlisting}

\subsection*{Dry run}

\dryrun{Grid $\begin{smallmatrix}0&2\\1&3\end{smallmatrix}$; start
$(0,0)$, end $(1,1)$ with id $3$.}

\begin{itemize}
  \item Elevation order: $0\,(0,0)$, $1\,(1,0)$, $2\,(0,1)$,
        $3\,(1,1)$.
  \item Activate $0$: no active neighbors. Activate $1\,(1,0)$: unite
        with $(0,0)$. Activate $2\,(0,1)$: unite with $(0,0)$.
  \item Activate $3\,(1,1)$: unite with active neighbors $(0,1)$ and
        $(1,0)$ -- now $(0,0)$ and $(1,1)$ are connected. Return
        elevation $\mathbf{3}$.
\end{itemize}

\subsection*{Complexity}

\begin{complexitybox}
\textbf{Time Complexity.} $O(n^2\log n)$ -- sorting the $n^2$ cells
dominates; each activation does $O(1)$ near-constant unions.
\textbf{Space Complexity.} $O(n^2)$ for the Disjoint Set, the cell
list, and the active grid.
\end{complexitybox}

\begin{takeawaybox}
When the question is ``the least threshold at which two points become
connected,'' add elements in threshold order and use Disjoint Set to
detect the moment of connection. This ``union over increasing time''
technique also solves problems phrased as minimizing the maximum edge on
a path -- for which a Dijkstra-style ``minimize the max'' (as in
Section~\ref{sec:graph-path-min-effort}) is the alternative.
\end{takeawaybox}
